 \documentclass[a4paper,openright,12pt]{report}
\usepackage{graphicx}
\usepackage[utf8]{inputenc}
\usepackage[hidelinks]{hyperref} 
\usepackage[hidelinks]{hyperref} 
\usepackage{multirow}
\usepackage{tikz}
\usepackage{fullpage}
\usepackage{parskip}
\usepackage{subfiles}
\usepackage{graphicx}
\usepackage{url}
\usepackage{xspace}
\usepackage{cite}\usepackage{cancel}
\usepackage{caption}
\usepackage{listings}
\usepackage{mathtools}
\usepackage{array}
\usepackage{listofsymbols} 
\usepackage[spanish, es-tabla, es-lcroman]{babel}
\usepackage{float}
\usepackage{longtable} 
\usepackage{nomencl}
\makenomenclature


\renewcommand{\nomname}{Glosario}
\newcommand*\circled[1]{\tikz[baseline=(char.base)]{
            \node[shape=circle,draw,inner sep=0.7pt] (char) {#1};}}
	
\hyphenation{op-tical net-works semi-conduc-tor cons-te-lla-tion co-nnecti-vi-ty}


\setlength{\parindent}{50pt}
\graphicspath{ {D:/Thesis/Escrito/img/} }
\newtheorem{q}{Question}[section]


\captionsetup{position=below}
\newtheorem{remarque}{Remarque}

\newcommand{\R}{\mathcal{R}}
\newcommand{\N}{\mathcal{N}}

\newcommand{\var}[1]{{\small\texttt{#1}} }
\newcommand{\lcavol}{\var{LCAVOL}}
\newcommand{\lweight}{\var{LWEIGHT}}
\newcommand{\age}{\var{AGE}}
\newcommand{\lbph}{\var{LBPH}}
\newcommand{\svi}{\var{SVI}}
% \newcommand{\lcp}{\var{LCP}}
\newcommand{\gleason}{\var{GLEASON}}
\newcommand{\pgg}{\var{PGG45}}
\newcommand{\lpsa}{\var{LPSA}}
\newcommand{\grad}{$^{\circ}$}
\usepackage{listings}
\setcounter{tocdepth}{4}
\setcounter{secnumdepth}{4}
\usepackage{xcolor}
\hypersetup{
    colorlinks,
    linkcolor={red!0!black},
    citecolor={black},
    urlcolor={blue!100!black}
}

\newcommand{\todo}[1]{\textcolor{blue}{TODO: #1}}

\usepackage[T1]{fontenc}   
\usepackage[utf8]{inputenc}

\fboxsep4mm

\fboxrule1.0pt
\usepackage{fancybox}
\usepackage{xcolor}

\definecolor{Zgris}{rgb}{0.87,0.85,0.85}

\newsavebox{\BBbox}
\newenvironment{DDbox}[1]{
\begin{lrbox}{\BBbox}\begin{minipage}{\linewidth}}
{\end{minipage}\end{lrbox}\colorbox{Zgris}{\usebox{\BBbox}} \\
[.5cm]}
\title{Proyecto Integrador Dardo Ariel Viñas Viscardi}
\begin{document}

%----------------------------------------------------------------------------------------
%	TITLE PAGE
%----------------------------------------------------------------------------------------

\begin{titlepage}
\begin{center}

\textsc{Universidad Nacional de Córdoba}\\[1.5cm] % University name

\includegraphics[width=0.3\textwidth]{img/logo_unc.pdf}~\\[1cm]

\textsc{\Large Proyecto Integrador Ingeniería en Computación }\\[0.5cm] % Thesis type

\rule{1.5\textwidth}{.4pt} % Horizontal line
\title{Title }\\[0.4cm] % Thesis title 
\textbf{\Huge{NetFi2}}
\rule{1.5\textwidth}{.4pt} %% Horizontal line
 
\begin{minipage}{0.4\textwidth}
\begin{flushleft} \large
\emph{Autor: Viñas Viscardi Dardo Ariel}\\


\end{flushleft}
\end{minipage}
\begin{minipage}{0.4\textwidth}
\begin{flushright} \large
\emph{Director: Dr. Ferreyra Pablo} \\
\emph{Codirector: Dr. Juan Fraire} \\
 
\end{flushright}
\end{minipage}\\[1cm]

\begin{table}[H]
\centering
\label{my-label}

\end{table}
 % University requirement text
%\textit{ }\\[0.4cm]

\large \textit{Laboratorio de Circuitos y Sistemas Robustos (LCSR)}\\ Facultad de Ciencias exactas Físicas y Naturales
\large \textit{}\\
\large \textit{}\\
\large \textit{}\\
\large \textit{ Lab. Techniques de l'Informatique et de la Microélectronique pour l'Architecture des systèmes intégrés (TIMA)}\\ Instituto Politécnico de Grenoble, Francia
 \large \textit{}\\
  \large \textit{}\\
   \large \textit{}\\
{\large \today}\\[2cm] % Date
%\includegraphics{Logo} % University/department logo - uncomment to place it
 
\vfill
\end{center}
\end{titlepage}


$\ $
\thispagestyle{empty} % para que no se numere esta pagina
% 


 % para comenzar la numeracion de paginas en números romanos

\chapter*{Resumen} % si no queremos que añada la palabra "capítulo"
\addcontentsline{toc}{chapter}{Resumen} % si queremos que aparezca en el índice
\markboth{RESUMEN}{RESUMEN} % encabezado
La inyección de falla es una técnica bien conocida para evaluar la sensibilidad de los circuitos integrados a los efectos de la radiación. En este trabajo, una metodología existente basada en emulación se extiende, actualiza y mejora bajo el nombre de \mbox{NETFI-3}. 

A diferencia de las soluciones anteriores, \mbox{NETFI-3} permite una eficiencia de recursos de hardware sin precedentes, flexibilidad de campaña, estimación altamente precisa de la tasa de error y lo más importante, una completa automatización de principio a fin de la metodología.
Esto se logra mediante un controlador de campaña incrustado dedicado a cargo de inyectar Single Event Upsets (SEU) y Single Event Transients (SETs) en bloques secuenciales y combinatorios con granularidad configurable. 

La metodología entera se ve automatizada, en partigular con el uso de pipelines para continou integration de GitLab, pero es simplemente una elección de implementación, ya que puede extenderse por completo a cualquier sistema actual de pipelines \( como también de forma local \). Se presentan extensos experimentos realizados en un estudio de caso de máquina bayesiana estocástica para demostrar las características y ventajas de \mbox{NETFI-3}.

%\chapter*{Introducción} % si no queremos que añada la palabra "capítulo"
%\addcontentsline{toc}{chapter}{Introducción} % si queremos que aparezca en el índice
%\markboth{Introducción}{Introducción} % encabezado
		


\cleardoublepage
\addcontentsline{toc}{chapter}{Lista de figuras} % para que aparezca en el índice de contenidos
% * <vinas.dardoariel@gmail.com> 2016-09-21T14:52:00.680Z:
%
% ^.
\listoffigures % indice de figuras

\cleardoublepage
\addcontentsline{toc}{chapter}{Lista de tablas} % para que aparezca en el indice de contenidos
\listoftables % indice de tablas

\tableofcontents % indice de contenidos
% \cleardoublepage

\addcontentsline{toc}{chapter}{Glosario}

\nomenclature{DUT}{Device Under \textit{test}} 
\nomenclature{FPGA}{Field-Programmable Gate Array } 
\nomenclature{TIMA}{Techniques de l'Informatique et de la Microélectronique pour l'Architecture des systèmes intégrés}
\nomenclature{SET}{Single Event Transient}
\nomenclature{SEU}{Single Event Upset}
\nomenclature{MBU}{Multi Bit Upset}
\nomenclature{MCU}{Multi Cell Upset}
\nomenclature{ESA}{European Space Agency}
\nomenclature{AXI}{Advanced eXtensible Interface}
\nomenclature{AMBA}{Advanced Microcontroller Bus Architecture}
\nomenclature{UART}{Universal Asynchronous Receiver Transmitter}
\nomenclature{VHDL}{VHSIC Hardware Description Language}
\nomenclature{RTL}{Register-Transfer Level}
\nomenclature{IDE}{Integrated Development Environment}
\nomenclature{JTAG}{Joint \textit{test} Action Group}
\nomenclature{LUT}{Lookup Table}
\nomenclature{FF}{Flip Flop}
\nomenclature{RIS}{Robust Integrated Systems}
\nomenclature{NETFI2}{Netlist Fault Injection 2}
\nomenclature{MODNET}{MODify NETlist}
\nomenclature{LET}{Lineal Energy Tranfer}
\nomenclature{NIEL}{Non Ionising Energy Loss}
\nomenclature{SEE}{Single Event Effect}
\nomenclature{SEL}{Single Event Latchup}
\nomenclature{EDIF }{Electronic Design Interchange Format}
\nomenclature{GeV}{Giga Electrón Volt}
\nomenclature{MeV}{Mega Electrón Volt}
\nomenclature{}{}


\printnomenclature

% \newpage


\pagenumbering{arabic} 

\subfile{content/chapters/01/chap01.tex}
\subfile{content/chapters/02/chap02.tex}
\subfile{content/chapters/03/chap03.tex}
\subfile{content/chapters/04/chap04.tex}
\subfile{content/chapters/05/chap05.tex}
\subfile{content/chapters/06/chap06.tex}



    
\bibliographystyle{content/bibtex/bst/unsrt}


\bibliography{content/bibtex/bib/bibliography}
Toda la bibliográfica consultada en sitios webs corresponde a consultas realizadas entre los meses Marzo y Octubre del año 2017. 

% \appendix

% \subfile{content/chapters/apend/01/apend01.tex}
% \subfile{content/chapters/apend/02/apend02.tex}
% \subfile{content/chapters/apend/03/apend03.tex}

\end{document}

\tableofcontents
\listoffigures
\listoftables
